\documentclass[a4paper, 11pt, titlepage]{article}
% \usepackage{geometry}
% \usepackage{longtable}
\usepackage{tabularx}
\usepackage{setspace}
\usepackage{bm}
\usepackage{amsmath}
\usepackage{hyperref}
\usepackage{graphicx,float}
\usepackage[export]{adjustbox}
\usepackage{listings}

\title{COMP5900C\\
Assignment 4\\
Reaction Diffusion}
\author{Gabriel Racz}

\begin{document}
\maketitle
\section{Implementation}
To facilitate real-time tuning of parameters, large output resolutions, and
interactive placement of chemicals I implemented the reaction diffusion
simulation system on the GPU using compute shaders. This is due to the uniform
computation performed accross each discrete grid cell, or pixel in our case. The
application is built using the Vulkan graphics API, but all of the
simulation logic is handled within the \texttt{reaction\_diffusion.glsl} and
\texttt{add\_chemicals.glsl} shaders.

We model the simulation domain using RGBA images. We represent each chemical
quantity as one color channel of the image. For each step of the simulation, we
apply our diffusion and reaction operations to the previous image state, and
save the results to our new image. This seperation of source and destination is
crucial for ensuring we can safely run our shader updates in parallel. Each
frame, we swap which image is the source and which is the destination. This
ensures that the source of each simulation step is the previous step's output.

The first step of each frame is to add/remove chemicals from the domain based on
user input. We support clearing all chemicals, globally seeding each pixel with a random
chemical quantity, clearing chemicals in a desired radius around the mouse
cursor, and adding chemicals in a radius around the mouse cursor. Exactly which
chemicals are added can be controlled through the GUI as specifying a desired
color (recall that each of R, G, and B represent a different chemical). This
interactivity allows designing patterns much better than purely offline
approaches.

To simulate the reaction diffusion process, 
\section{Results}
\end{document}